\section{Sustainability and Frontier-Model Dependence}
\label{sec:sustainability}

The pattern is not free. We owe the reader an explicit account of
what F(AI)\textsuperscript{2}R consumes and what it depends on.

\paragraph{Resource cost.}
Every authoring pass and every audit pass burns compute, and the
verification ladder amplifies this: a contested claim is read by an
agent, by a human, sometimes by a second agent during readability
review, and again by the FAIR aligner. A heavily-revised paper can
drive dozens of model invocations per claim. The energy and water
footprint of frontier-model inference is non-trivial and is paid by
someone, somewhere, regardless of who pays for the API call. The
seventh integrated practice (\S\ref{sec:eight}, item~7) is silent on
this; we add a request here that disclosure under that practice
\emph{include the order-of-magnitude inference budget consumed by the
manuscript}, declared at submission time alongside the rung
distribution.

\paragraph{Dependence on frontier models.}
The pattern's effective behaviour depends on the capability of the
underlying model. As frontier models change --- through vendor
turnover, version deprecation, capability shifts, training-data
updates, or refusal-policy changes --- the pattern's outputs will
shift with them. The transcripts of Practice~1 and the
tool-and-model version manifest are partial mitigation: a future
researcher can at least know which model produced which sentence,
even if they can no longer invoke it. The receipts are preserved;
the capability is not portable. We treat this as a stranded-asset
risk and recommend that any \texttt{fair2r:Activity} record carry an
explicit \texttt{fair2r:modelId} property pinned to a stable
identifier (model name, version, and where available the model card
URI), so that a deprecated agent at minimum has a known headstone.

\paragraph{Equity.}
Research groups without access to frontier models cannot apply this
pattern at the same effective rung; the verification ladder requires
that an AI agent at \textsc{ai-confirmed} actually understands the
source it has read, and the floor for ``understands'' is set by the
state of the art. F(AI)\textsuperscript{2}R \emph{describes} the
work; it does not \emph{equalise} access to it. The pattern is
therefore not a neutral standard. We name this honestly and offer
two partial mitigations: (i) smaller-model fallbacks for routine
audit passes (layout, hedging-density, list-of-lists detection) that
do not need frontier-grade reasoning; (ii) the structural property
that an honestly-declared lower-rung distribution at submission time
is preferable to a forged higher-rung distribution from a richer
group, and the verification ladder makes the difference legible.

\paragraph{Long-term archivability.}
A paper whose claims depend on a now-deprecated model carries a
replay risk: the artefacts can still be read but the agents cannot
be re-executed against new sources. The mitigation, again, is
preservation: transcripts, prompts, the version manifest, and the
provenance graph survive model deprecation in a way that the
inference itself does not. We recommend that the
F(AI)\textsuperscript{2}R cell mapping for R1.2 (provenance) include
a hash of the prompt set and a recorded model identifier, so that an
archive can preserve the recipe even when the kitchen is gone.

The sustainability story is not optional. A practice that produces
auditable artefacts at the cost of disposable inference is, on a
long enough horizon, neither sustainable nor auditable. Naming the
cost is the price of admission.
